\section{Guided Solutions and Commentary}
\chapterline

A polished mathematical proof is often intimidating because it conceals all the scaffolding used to construct it. When you read a finished proof in a textbook, the author seems to magically know the exact $\delta$, the perfect algebraic substitution, or the clever contradiction hypothesis on line one.

In reality, mathematical discovery runs in reverse. You stumble, draw diagrams, test small numbers, write backwards equations on scrap paper, and hit dead ends. Only after the territory is thoroughly explored do you turn around and write down the clean, forward deductive path.

This chapter is designed to make that hidden scaffolding visible. For each problem, we present both tracks of mathematical thought: the exploratory search and the polished deduction.

\begin{historybox}
The contrast between how mathematics is discovered and how it is published is illustrated by two of the greatest figures in history: Leonhard Euler and Carl Friedrich Gauss.

Euler was remarkably candid. In his hundreds of memoirs, he invited the reader directly into his workshop. In solving the famous Basel problem in 1735 (finding $\sum_{n=1}^\infty \frac{1}{n^2} = \frac{\pi^2}{6}$), Euler calculated sums to twenty decimal places, factored infinite sine polynomials as if they were simple quadratics, and transparently documented his intuitive leaps long before Weierstrass developed the rigorous theory of infinite product expansions.

Gauss took the opposite approach. Gauss was obsessed with perfection, adopting the motto \emph{Pauca sed matura} (``Few, but ripe''). He famously stated: ``When a fine building is completed, the scaffolding should no longer be visible.'' Gauss systematically erased all scratchwork, false starts, and intuitive heuristic bridges from his published papers, leaving behind monolithic, flawless deductive gems that stunned and intimidated his contemporaries. In this chapter, we tear down Gauss's finished facade to reveal Euler's working scaffolding.
\end{historybox}

\subsection{The Architecture of Proof Construction}

Before examining individual problems, consider how the two tracks interact during the birth of any proof:

\begin{center}
\begin{tikzpicture}[scale=1.05, >=Stealth]
  \tikzset{
    tracklane/.style={very thick, rounded corners=4pt},
    stepbox/.style={thick, rounded corners=3pt, inner sep=5pt, font=\small},
    arrowstyle/.style={->, thick, shorten >=2pt, shorten <=2pt},
    bridgearrow/.style={->, very thick, draw=structurepurple!90!black, dashed}
  }

  % Headers
  \node[font=\normalsize\scshape, text=quotientorange!90!black] at (2.4, 5.4) {Discovery Track (Scratchpad)};
  \node[font=\normalsize\scshape, text=targetgreen!90!black] at (8.0, 5.4) {Presentation Track (Formal Proof)};

  % Lanes
  \draw[tracklane, draw=quotientorange!80!black, fill=quotientorange!5] (0.2, -0.4) rectangle (4.6, 4.9);
  \draw[tracklane, draw=targetgreen!80!black, fill=targetgreen!5] (5.8, -0.4) rectangle (10.2, 4.9);

  % Steps: Discovery
  \node[stepbox, draw=quotientorange!80!black, fill=quotientorange!15] (d1) at (2.4, 4.2) {1. Target: Look at Conclusion};
  \node[stepbox, draw=quotientorange!80!black, fill=quotientorange!15] (d2) at (2.4, 2.9) {2. Work Backwards from Goal};
  \node[stepbox, draw=quotientorange!80!black, fill=quotientorange!15] (d3) at (2.4, 1.6) {3. Unpack Hidden Constraints};
  \node[stepbox, draw=quotientorange!80!black, fill=quotientorange!15] (d4) at (2.4, 0.3) {4. Discover Candidate $N, \delta, a$};

  \draw[arrowstyle, draw=quotientorange!90!black] (d1) -- (d2);
  \draw[arrowstyle, draw=quotientorange!90!black] (d2) -- (d3);
  \draw[arrowstyle, draw=quotientorange!90!black] (d3) -- (d4);

  % Steps: Presentation
  \node[stepbox, draw=targetgreen!80!black, fill=targetgreen!15] (p1) at (8.0, 4.2) {1. Declare Preconditions / Let $x$};
  \node[stepbox, draw=targetgreen!80!black, fill=targetgreen!15] (p2) at (8.0, 2.9) {2. Introduce Discovered Candidate};
  \node[stepbox, draw=targetgreen!80!black, fill=targetgreen!15] (p3) at (8.0, 1.6) {3. Forward Deductive Chain};
  \node[stepbox, draw=targetgreen!80!black, fill=targetgreen!15] (p4) at (8.0, 0.3) {4. Conclude Target Satisfied \checkmark};

  \draw[arrowstyle, draw=targetgreen!90!black] (p1) -- (p2);
  \draw[arrowstyle, draw=targetgreen!90!black] (p2) -- (p3);
  \draw[arrowstyle, draw=targetgreen!90!black] (p3) -- (p4);

  % Cross-track bridge
  \draw[bridgearrow] (4.4, 0.3) to[bend right=15] node[below, font=\scriptsize\bfseries, text=structurepurple!90!black] {Candidate feeds Step 2} (6.0, 2.9);
\end{tikzpicture}
\end{center}

\subsection{Controversies, False Starts, and the Genius Myth}

The presentation tradition in mathematics has created enduring pedagogical damage and widespread misconceptions:

\begin{itemize}
  \item \textbf{The Myth of the Effortless Genius}: Because finished textbooks read like continuous cascades of brilliant deductions, students often conclude that real mathematicians never make mistakes or explore dead ends. This produces crippling imposter syndrome. In truth, every mathematician spends 90\% of their time in the messy, disoriented scratchpad phase. A proof is not a recording of the thinking process; it is a clean legal argument written after the investigation has succeeded.
  \item \textbf{The Backward Deduction Fallacy}: The single most dangerous beginner error is writing backwards scratchwork directly into the final proof. If you write:
  \[
  |f(x) - L| < \eps \impliedby |3x - 6| < \eps \impliedby 3|x - 2| < \eps \impliedby |x - 2| < \frac{\eps}{3},
  \]
  and end with ``Choose $\delta = \eps/3$,'' you have reversed the direction of logical flow. In a forward proof, the hypothesis $|x - 2| < \delta$ must be the starting premise that drives forward to the conclusion $|f(x) - L| < \eps$.
  \item \textbf{Premature Formalization}: Attempting to write formal symbolic lines before developing an intuitive mental model or sketching a counterexample inevitably leads to empty symbol pushing. If you cannot explain the reason for a claim in plain words to a friend, no amount of quantifiers will save the proof.
\end{itemize}

\subsection{Problem 1: Direct Proof and Parity}

\begin{theorem}
The sum of two odd integers is an even integer.
\end{theorem}

\paragraph{The Mental Model \& False Starts.}
A common beginner impulse is to test examples: $3 + 5 = 8$, $7 + 11 = 18$. Examples build confidence, but they cannot prove an infinite claim. Another common blunder is to write $a = 2k + 1$ and $b = 2k + 1$, accidentally forcing $a$ and $b$ to be the identical number! Different arbitrary objects must receive independent variable names.

\paragraph{The Scratchpad.}
We need $a + b$ to look like $2 \cdot (\text{integer})$. Since $a$ is odd, $a = 2m + 1$. Since $b$ is odd, $b = 2n + 1$ for some integers $m, n$. Adding them gives $2m + 2n + 2 = 2(m + n + 1)$. Because integers are closed under addition, $m + n + 1$ is an integer. The bridge is complete.

\begin{proof}
Let $a$ and $b$ be arbitrary odd integers. By definition of odd integers, there exist integers $m, n \in \Z$ such that
\[
a = 2m + 1 \quad \text{and} \quad b = 2n + 1.
\]
We compute the sum:
\begin{align*}
a + b &= (2m + 1) + (2n + 1) \\
&= 2m + 2n + 2 \\
&= 2(m + n + 1).
\end{align*}
Let $k = m + n + 1$. Since $m, n, 1 \in \Z$ and the integers are closed under addition, $k \in \Z$. Thus $a + b = 2k$ with $k \in \Z$, which satisfies the definition of an even integer. Since $a$ and $b$ were arbitrary, the sum of any two odd integers is even.
\end{proof}

\subsection{Problem 2: Unwinding Quantifiers in Surjections}

\begin{theorem}
If $f: A \to B$ and $g: B \to C$ are surjective functions, then the composition $g \circ f: A \to C$ is surjective.
\end{theorem}

\paragraph{The Mental Model \& False Starts.}
A beginner often starts by picking an element $a \in A$ and pushing forward: $f(a) \in B$, $g(f(a)) \in C$. But surjectivity is a promise about the \emph{target}, not the domain! The definition states: for every $c \in C$, there exists an input mapping to it. Starting at $A$ is moving in the wrong direction.

\paragraph{The Scratchpad.}
We must start with an arbitrary $c \in C$. We need to locate some $a \in A$ with $(g \circ f)(a) = c$.
\begin{itemize}
  \item What map touches $C$? The function $g: B \to C$. Because $g$ is surjective, there is some $b \in B$ with $g(b) = c$.
  \item Now we have $b \in B$. What map touches $B$? The function $f: A \to B$. Because $f$ is surjective, there is some $a \in A$ with $f(a) = b$.
  \item Putting them together: $(g \circ f)(a) = g(f(a)) = g(b) = c$. We found our $a$.
\end{itemize}

\begin{proof}
Let $c \in C$ be arbitrary. We must exhibit an element $a \in A$ such that $(g \circ f)(a) = c$.

Because $g: B \to C$ is surjective, there exists an element $b \in B$ such that $g(b) = c$. 

Next, because $f: A \to B$ is surjective and $b \in B$, there exists an element $a \in A$ such that $f(a) = b$.

Applying the definition of function composition:
\[
(g \circ f)(a) = g(f(a)) = g(b) = c.
\]
Since $c \in C$ was arbitrary, we have demonstrated that every element of $C$ possesses a preimage in $A$ under $g \circ f$. Therefore $g \circ f$ is surjective.
\end{proof}

\subsection{Problem 3: Type Discipline in Power Sets}

\begin{theorem}
For any sets $A$ and $B$, if $A \subseteq B$, then $\pow(A) \subseteq \pow(B)$.
\end{theorem}

\paragraph{The Mental Model \& False Starts.}
This is a classic type-confusion trap. Students frequently write: ``Let $x \in A$. Since $A \subseteq B$, $x \in B$. Therefore $\pow(A) \subseteq \pow(B)$.'' This is invalid because it proves $A \subseteq B$ (which was already given), not $\pow(A) \subseteq \pow(B)$! 

An element of a power set $\pow(A)$ is not a point; it is an \emph{entire subset} $S \subseteq A$.

\begin{proof}
Assume $A \subseteq B$. To prove that $\pow(A) \subseteq \pow(B)$, let $S \in \pow(A)$ be arbitrary.

By definition of the power set, $S \in \pow(A)$ means that $S \subseteq A$.

We must show that $S \in \pow(B)$, which means showing that $S \subseteq B$.

Let $x \in S$ be arbitrary. Since $S \subseteq A$, we have $x \in A$. Since $A \subseteq B$, we have $x \in B$. 

Because every element of $S$ is an element of $B$, it follows that $S \subseteq B$.

By definition of the power set, $S \subseteq B$ implies $S \in \pow(B)$. Since $S \in \pow(A)$ was arbitrary, we conclude that $\pow(A) \subseteq \pow(B)$.
\end{proof}

\subsection{Problem 4: Proof by Contradiction and Density}

\begin{theorem}
There is no smallest positive real number.
\end{theorem}

\paragraph{The Mental Model.}
If someone claims that a mountain has a lowest peak above sea level, you produce a point halfway between that peak and sea level. In $\R$, between any positive number and zero lies another positive number.

\begin{proof}
Suppose, for contradiction, that there exists a smallest positive real number. Call this number $r$.

By assumption:
\begin{enumerate}
  \item $r > 0$, and
  \item For every real number $x > 0$, we have $r \le x$.
\end{enumerate}
Now consider the number $s = \frac{r}{2}$. 

Since $r > 0$, dividing by $2$ preserves the inequality, so $s = \frac{r}{2} > 0$. Thus $s$ is a positive real number.

However, since $r > 0$, we have $\frac{r}{2} < r$, which means $s < r$.

This directly contradicts condition (2), which asserted that $r$ is less than or equal to \emph{every} positive real number. 

Because our assumption led to an explicit logical contradiction, that assumption must be false. Therefore, there is no smallest positive real number.
\end{proof}

\subsection{What Guided Scaffolding Unlocks---and Where It Hits a Wall}

Making the two tracks of proof visible unlocks authentic mathematical independence:
\begin{enumerate}
  \item \textbf{Proof Debugging}: You can diagnose exactly where a proof broke down: did you conflate element types? Did you push from the domain instead of pulling from the target? Did you reverse an implication?
  \item \textbf{Liberation from Memory}: You no longer need to memorize hundreds of proof recipes. By unpacking definitions and working backwards from the conclusion, the structure of the proof reveals itself naturally.
  \item \textbf{Auditable Rigor}: By separating exploratory intuition from forward deduction, your arguments become clear and convincing to any reader.
\end{enumerate}

\paragraph{Honest Downsides and Limitations.}
Yet guided commentary has natural limitations:
\begin{itemize}
  \item \textbf{Templates Cannot Solve Novel Research}: Following standard proof architectures works reliably for textbook theorems. But in open research problems, standard templates often fail. Progress requires radical creative leaps---introducing an auxiliary concept (like homology or schemes) that seems totally unrelated to the original problem.
  \item \textbf{Deductive Rules Cannot Generate Conjectures}: Logic verifies truth; it does not discover it. As Henri Poincar\'e observed, logic is the instrument of proof, but intuition is the instrument of discovery.
\end{itemize}

\paragraph{Where We Go Next in the Book.}
You have seen the scaffolding, analyzed the historical controversies, and mastered the grammar of pure mathematics. It is time to test your fluency. In Chapter~41 (\emph{Final Exam: From Beginner to Ready}), we present the capstone assessment unifying all five foundational pillars of pure mathematics before you launch into advanced mathematical study.

\newpage
